% from ega4-17-fr.tex
\oldpage[IV-4]{66}
est un $B$-module de type fini (16.3.9), donc \emph{séparé} pour la topologie $\mathfrak n$-préadique (où $\mathfrak n$ est l'idéal maximal de $B$) ($0_{\mathrm I}$, 7.3.5); il revient par suite au même de dire que $\Omega^1_{B/A}=0$ ou que $\widehat{\Omega}^1_{B/A}=0$; donc (0, 20.7.4), il revient au même de dire que $B$ est une $A$-Algèbre formellement non ramifiée pour les topologies \emph{discrètes}, ou pour les topologies \emph{préadiques}. Ceci prouve l'équivalence des conditions e) et f'). Si $\mathfrak m$ est l'idéal maximal de $A$, on a
\[
k=A/\mathfrak m=\widehat A/\mathfrak m\widehat A,
\qquad
\widehat B\otimes_{\widehat A}k=\widehat B/\mathfrak m\widehat B
=\widehat B\otimes_B(B/\mathfrak m B),
\]
et par suite ($0_{\mathrm I}$, 7.3.5) $\widehat B/\mathfrak m\widehat B$ est le complété de $B/\mathfrak m B=B\otimes_A k$ pour la topologie $\mathfrak n$-préadique; ceci prouve l'équivalence de d'') et de f). Enfin, lorsque $k(x)=k(y)$ ou lorsque $k$ est séparablement clos, la condition f) entraîne que l'homomorphisme $\widehat A/\mathfrak m\widehat A\to\widehat B/\mathfrak m\widehat B$ est bijectif; la condition f) implique d'autre part que $\widehat B$ est une $\widehat A$-algèbre \emph{quasi-finie} ($0_{\mathrm I}$, 7.4.4), donc \emph{finie} puisque $\widehat A$ est complet et $\widehat B$ séparé pour la topologie $\mathfrak m$-préadique, $\mathfrak m\widehat B$ étant un idéal de définition de $\widehat B$ ($0_{\mathrm I}$, 7.4.1). L'homomorphisme $\widehat A\to\widehat B$ est donc surjectif en vertu du lemme de Nakayama. Donc f) entraîne alors f''), et la réciproque est évidente.

\begin{env}[17.4.5]
\EGAReaderTarget{ega4.c04.statement.17-4-5}{IV.17.4.5}
Étant donnés un $S$-préschéma $Y$ et deux $S$-morphismes $f:X\to Y$, $g:X\to Y$, on en déduit canoniquement un $S$-morphisme $(f,g)_S:X\to Y\times_S Y$. Nous appellerons \emph{préschéma des coïncidences} de $f$ et $g$ l'image réciproque par $(f,g)_S$ de la diagonale $\Delta_{Y|S}$; c'est donc un sous-préschéma de $X$, qui est \emph{fermé} lorsque $Y$ est un $S$-schéma (I, 5.4.1).
\end{env}

\begin{proposition}[17.4.6]
\label{IV.17.4.6}
Soient $h:Y\to S$ un morphisme non ramifiée, et soient $f:X\to Y$, $g:X\to Y$ deux $S$-morphismes. Alors le préschéma $C$ des coïncidences de $f$ et $g$ est un sous-préschéma induit sur un ouvert de $X$; si de plus $Y$ est un $S$-schéma (I, 5.4.1), $C$ est un sous-préschéma fermé de $X$.
\end{proposition}

En effet, puisque $\Delta_h:Y\to Y\times_S Y$ est une immersion ouverte (17.4.2), l'image réciproque par $(f,g)_S$ de $\Delta_{Y|S}$ est un sous-préschéma induit sur un ouvert de $X$ (I, 4.4.1). La dernière assertion résulte de (17.4.5).

\begin{corollary}[17.4.7]
\label{IV.17.4.7}
Sous les hypothèses de (17.4.6), soit $x$ un point de $X$ tel que les deux morphismes composés $\Spec(k(x))\to X\xrightarrow{f}Y$ et $\Spec(k(x))\to X\xrightarrow{g}Y$ soient égaux. Alors il existe un voisinage ouvert $U$ de $x$ tel que $f|U=g|U$. Si de plus $Y$ est un $S$-schéma, il existe un voisinage ouvert et fermé $X'$ de $x$ dans $X$ tel que $f|X'=g|X'$. Si enfin on suppose en outre que $X$ est connexe, on a $f=g$.
\end{corollary}

Cela résulte de (17.4.6) et de (I, 5.3.17).

\begin{corollary}[17.4.8]
\label{IV.17.4.8}
Sous les hypothèses de (17.4.6), supposons que le morphisme structural $\varphi=h\circ f=h\circ g$ de $X$ dans $S$ soit fermé. Soit $s$ un point de $S$; soit $X_s$ le $k(s)$-préschéma $\varphi^{-1}(s)$ et supposons que les deux morphismes composés $X_s\to X\xrightarrow{f}Y$ et $X_s\to X\xrightarrow{g}Y$ soient égaux. Alors il existe un voisinage ouvert $V$ de $s$ dans $S$ tel que $f|\varphi^{-1}(V)=g|\varphi^{-1}(V)$. Si de plus $Y$ est un $S$-schéma et si $\varphi$ est ouvert, on peut prendre $V$ ouvert et fermé. Si enfin on suppose de plus $S$ connexe, on a $f=g$.
\end{corollary}

Il résulte de (17.4.7) que le préschéma $C$ des coïncidences de $f$ et $g$ est induit sur un ouvert de $X$ et contient $X_s$. Comme $\varphi$ est fermé, il existe un voisinage ouvert $V$ de $s$ tel que $\varphi^{-1}(V)\subset C$. Si de plus $Y$ est un $S$-schéma, $C$ est fermé, donc $\varphi(X-C)$ est à la

